\section{A finite Sobolev theorem and an explicit extension contract}
\label{sec:sobolev}

The article relates Euclidean isoperimetry to the sharp \(W^{1,1}\)
Sobolev inequality. There is a useful exact constructive version before
any infinite-dimensional function space is introduced.

Fix \(n\geq2\) and put \(p=n/(n-1)\). Let
\(E_1\supseteq\cdots\supseteq E_m\) be nested regions with supplied
volumes \(v_j\), perimeters \(P_j\), and presentations to which
Theorem~\ref{thm:euclidean-isoperimetry} applies. Zero-volume entries
can be omitted when they are decidably empty; alternatively their
individual inequality can be supplied directly. Let
\(0=t_0<t_1<\cdots<t_m\) be rational and
\(\delta_j=t_j-t_{j-1}\). The associated nonnegative simple function is
\(u=\sum_{j=1}^m\delta_j\mathbf1_{E_j}\). Its entire numerical data
consist of the finitely many shell volumes
\(w_i=v_i-v_{i+1}\geq0\), where \(v_{m+1}=0\):
\[
 \|u\|_p=\left(\sum_{i=1}^m w_i t_i^p\right)^{1/p},
 \qquad J(u)=\sum_{j=1}^m\delta_jP_j.
\]
Fractional powers are the primitive recursively named roots described
earlier. Neither an integral over all real points nor a completed
Lebesgue function space is needed to define these quantities.

\begin{lemma}[Finite weighted triangle inequality]
For nonnegative weights \(w_i\), the quantity
\(\|a\|_p=(\sum_i w_i|a_i|^p)^{1/p}\) satisfies
\(\|a+b\|_p\leq\|a\|_p+\|b\|_p\).
\end{lemma}
\begin{proof}
We give the finite proof to expose its strength. Arithmetic--geometric
mean on \(n\) numbers, \(n-1\) copies of \(A^p\) and one copy of
\(B^n\), gives Young's inequality
\[
 AB\leq\frac{n-1}{n}A^p+\frac1nB^n.
\]
After normalization and summation this gives the weighted finite
H\"older inequality with conjugate exponents \(p,n\):
\[
 \sum_iw_i A_i B_i
 \leq\left(\sum_iw_i A_i^p\right)^{1/p}
       \left(\sum_iw_i B_i^n\right)^{1/n}.
\]
To avoid any undecidable division by zero, normalize using
\(\|A\|_p+\varepsilon\) and \(\|B\|_n+\varepsilon\) for positive
rational \(\varepsilon\), and then use the order definition as
\(\varepsilon\) runs through prescribed dyadics. The normalized sums
are at most one, so the same calculation proves the desired bound with
the two extra \(\varepsilon\)'s; there is no existence search.

Apply H\"older to \(A_i=|a_i|\), and then to \(A_i=|b_i|\),
with \(B_i=|a_i+b_i|^{p-1}\). Since \((p-1)n=p\),
\[
 \|a+b\|_p^p
 \leq(\|a\|_p+\|b\|_p)\|a+b\|_p^{p-1}.
\]
This implies the triangle inequality constructively: if its conclusion
failed by a positive rational margin, then \(\|a+b\|_p\) would have
an explicit positive lower bound, permitting division and contradicting
that margin. Equivalently perform the comparison on rational root
intervals at a precision fixed by the margin. This proves the order
inequality without real trichotomy. Iteration gives the finite-sum form.
\end{proof}

\begin{theorem}[Finite layered sharp Sobolev inequality]
\label{thm:layered-sobolev}
For the supplied nested data,
\[
 \left(\sum_{i=1}^m w_i t_i^{n/(n-1)}\right)^{(n-1)/n}
 \leq \frac1{n\omega_n^{1/n}}\sum_{j=1}^m\delta_jP_j.
\]
\end{theorem}
\begin{proof}
View the function as a vector on the \(m\) shells. Its \(j\)-th
layer is the vector equal to \(\delta_j\) on shells with index
\(i\geq j\), and zero on the other shells. Its weighted \(p\)-norm
is \(\delta_jv_j^{1/p}\). The finite triangle inequality gives
\(\|u\|_p\leq\sum_j\delta_jv_j^{(n-1)/n}\).
Apply Euclidean isoperimetry separately to each \(E_j\), then sum.
All sums have length at most \(m\), and positivity of \(\omega_n\)
was proved by an explicit contained cube.
\end{proof}

\subsection{Exact coarea for rational grid functions}

Consider finitely many rational grid cells with rational nonnegative
values, extended by zero outside a bounding rectangle. List the distinct
positive values as \(t_1<\cdots<t_m\). The superlevel sets
\(E_j=\{u\geq t_j\}\) are finite unions of grid cells. Let
\(P(E_j)\) be the sum of the areas of their exposed faces.
The finite box-union calculation in Section~\ref{sec:euclidean} proves
that these are valid outer-perimeter bounds, so the preceding theorem
applies unconditionally to this finite input.

For each interface face \(F\), let \(u_F^+,u_F^-\) be the values
on its two sides, one of which is zero on an exterior face. The exact
coarea identity is
\[
 \sum_j\delta_jP(E_j)
 =\sum_F\operatorname{area}_{n-1}(F)|u_F^+-u_F^-|.
\]
Indeed, if the two adjacent values are \(a<b\), a fixed face contributes
precisely for indices satisfying \(a<t_j\leq b\); the corresponding
\(\delta_j\)'s telescope to \(b-a\). Summing this
identity over the finite list of faces proves the formula. This gives a
finite BV form of Sobolev's inequality with the Euclidean isoperimetric
constant and an explicitly
computable right-hand side. It is not being confused with a classical
gradient of a discontinuous grid function.

\subsection{What must be supplied for a \(W^{1,1}\) interpretation}

A precise extension is as follows. Supply a named target norm \(N\),
a named gradient integral \(G\), and a primitive recursive sequence of
layered data \(u_k\) satisfying, with proofs,
\[
 N\leq\|u_k\|_p+2^{-k},\qquad J(u_k)\leq G+2^{-k}.
\]
Then Theorem~\ref{thm:layered-sobolev} gives
\[
 N\leq \frac{G}{n\omega_n^{1/n}}
       +\left(1+\frac1{n\omega_n^{1/n}}\right)2^{-k}
\]
for every \(k\), hence \(N\leq G/(n\omega_n^{1/n})\).
This proves the analytic Sobolev inequality for any presentation that
supplies these norm and coarea approximation data. Constructing such data
for every abstract classical \(W^{1,1}\) element is not a theorem claimed
here. In particular, naive approximation by axis-aligned step functions
does not automatically preserve the Euclidean gradient norm: it may
converge to an anisotropic variation instead. The explicit condition on
\(J(u_k)\) prevents that false inference.

The constant cannot be decreased for the finite layered class: the single
layer \(u=\mathbf1_{B_R}\) has
\(\|u\|_p=(\omega_nR^n)^{(n-1)/n}\) and
\(J(u)=n\omega_nR^{n-1}\), which give equality. These ball data are
effective by the preceding section. The Sobolev sharpness approximation
can also be made quantitative with a radial Lipschitz linear transition on the
shell \(R\leq|x|\leq R+\varepsilon\): its gradient integral is
\(\omega_n((R+\varepsilon)^n-R^n)/\varepsilon\), and the norm error
is bounded by the volume of that shell to the power \(1/p\). The
binomial estimates already supplied give all rates. This verifies
sharpness by an explicit approximation family, without asserting a
universal effective coarea theorem.

\section{Spherical statements: caps, finite cap unions, and complements}
\label{sec:spherical}

The general spherical isoperimetric inequality for every presented
simple closed curve is not proved in this report. The following
substantial explicitly represented subclass, together with the exact
equality examples and complement symmetry, is completely elementary.
This scope distinction matters: a calculation for latitude circles is
not a proof that all spherical curves have greater perimeter.

Let the sphere have radius \(R>0\), with a supplied positive lower bound,
and let a cap have angular radius \(\theta\in[0,\pi]\), named with
primitive recursive approximations. The standard parametrization is
\[
 (\phi,t)\longmapsto
 (R\sin\phi\cos t,R\sin\phi\sin t,R\cos\phi),
 \qquad 0\leq\phi\leq\theta,\quad0\leq t\leq2\pi.
\]
Its boundary speed is \(R\sin\theta\), and the norm of the cross
product of its two partial derivatives is \(R^2\sin\phi\).
These are finite trigonometric identities; the area and length integrals
have explicit antiderivatives and can equally be defined by rational
Riemann sums with the elementary sine and cosine moduli. Thus
\[
 A=2\pi R^2(1-\cos\theta),\qquad
 L=2\pi R\sin\theta.
\]
Using \(\sin^2\theta+\cos^2\theta=1\), direct substitution gives
\begin{equation}\label{eq:cap-equality}
 L^2=4\pi A-\frac{A^2}{R^2}.
\end{equation}
The polar parametrization's degeneracy at a pole is harmless: truncating
a polar strip changes area by its explicit cap-area formula, tending to
zero with a prescribed modulus. No division by \(\sin\phi\) is used.

\begin{theorem}[Spherical isoperimetry for disjoint cap unions]
Let a finite collection of caps have disjoint interiors, with supplied
evidence that their boundary lengths add (for example, boundaries are
disjoint or meet at finitely many points). Let \(A\) be the sum of their
areas and \(L\) the sum of their boundary lengths. Then
\[
 L^2\geq4\pi A-\frac{A^2}{R^2}.
\]
The same inequality holds for the complement of the union on the sphere.
\end{theorem}
\begin{proof}
Write the cap data as \(A_i,L_i\) and use
\eqref{eq:cap-equality}. Expansion gives the exact identity
\[
 L^2-4\pi A+\frac{A^2}{R^2}
 =2\sum_{i<j}L_iL_j+\frac2{R^2}\sum_{i<j}A_iA_j\geq0.
\]
For the complement, the area is \(A'=4\pi R^2-A\) and its boundary
length is \(L\). Direct algebra shows
\(4\pi A'-(A')^2/R^2=4\pi A-A^2/R^2\).
This proves the complement assertion without selecting which side of a
simple closed curve is the smaller side.
\end{proof}

If there are two caps with explicitly positive areas, the displayed
deficit is explicitly positive. A single nondegenerate cap gives equality.
For \(R=1\), the right side is \(A(4\pi-A)\), exactly the spherical
formula in the article. Scaling lengths by \(R\) and areas by \(R^2\)
explains the general-radius formula. The theorem proves a finite-union
extension as well as the equality examples, while leaving general
spherical rearrangement and general spherical rigidity outside its scope.

